% ============================================================
%  CHAPTER 7 — Sprint 5: Conversational Assistant
%  Reusable content block (no \documentclass, no preamble).
%  Author: Sarra Mediouni
% ============================================================
\chapter{Sprint 5: Conversational Assistant}

\section{Introduction}

By the end of Sprint~4, the gaming centre had a star schema, a four-page
Power~BI report, and a full layer of forecasting, segmentation, and
anomaly detection. The numbers were all there. What was still missing was
a way for the owner to actually \emph{ask} a question without having to
remember which dashboard page hides which answer.

Sprint~5 closes that gap with a conversational assistant. The idea is
simple: the owner types (or speaks) a question in plain language --- in
English, French, or Arabic --- and the system replies with a written
answer and, when relevant, a chart. The assistant sits on top of the same
CSV outputs Power~BI consumes, so there is one source of truth.

\section{Sprint 5 Backlog}

Table~\ref{tab:sprint5_backlog} lists the tasks committed to this sprint.

\begin{longtable}{>{\bfseries}p{1.3cm} p{4.4cm} p{7cm}}
  \caption{Sprint 5 Backlog.\label{tab:sprint5_backlog}} \\
  \toprule
  \textbf{ID} & \textbf{User Story} & \textbf{Tasks} \\
  \midrule
  \endfirsthead
  \toprule
  \textbf{ID} & \textbf{User Story} & \textbf{Tasks} \\
  \midrule
  \endhead
  \bottomrule
  \endfoot

  US-18 & As an owner, I want to ask questions in plain language and get
          answers with charts.
        & Build a Streamlit front-end (\texttt{app.py}) wired to a Gemini-backed
          agent (\texttt{chatbot/agent.py}); load the star-schema and forecast
          CSVs as a shared \texttt{DataContext} (\texttt{chatbot/data\_loader.py}). \\[4pt]

  US-19 & As an owner, I want to use the assistant in English, French, or
          Arabic.
        & Translate the UI string bundle in \texttt{app.py} (\texttt{STRINGS}
          dict, three locales); pass a language instruction to the system
          prompt so the model also answers in the chosen language. \\[4pt]

  US-20 & As an owner, I want to talk to the assistant instead of typing.
        & Add a custom Streamlit component (\texttt{chatbot/voice\_component.py})
          backed by an HTML/JS frontend (\texttt{chatbot/voice\_html/}) that
          uses the browser's Web Speech API; pass the selected language as
          the recognition locale (\texttt{en-US}, \texttt{fr-FR}, or \texttt{ar-TN}). \\[4pt]

  US-21 & As an owner, I want a chart drawn on demand when one would help.
        & Implement \texttt{chatbot/chart\_generator.py} on top of Plotly
          Express; have the agent emit a structured \texttt{chart\_spec}
          JSON block when a chart is useful and render it in the right
          panel of the layout. \\
\end{longtable}

\section{Architecture Overview}

Sprint~5 adds a fifth layer on top of the four-layer pipeline introduced
in Sprint~0. Figure~\ref{fig:chatbot_arch} sketches the flow at runtime.

\begin{figure}[H]
  \centering
  % *TODO* Render a sequence/arch diagram for the chatbot. Suggested boxes:
  % User -> Streamlit (app.py) -> [DataContext (CSVs)] + [voice_component]
  %       -> agent.ask() -> Gemini API -> agent._parse_response() -> chart_generator
  %       -> Plotly figure + answer + suggestions -> back to Streamlit.
  \includegraphics[width=0.92\textwidth]{images/chatbot_arch.png}
  \caption{Runtime architecture of the conversational assistant.}
  \label{fig:chatbot_arch}
\end{figure}

The pieces are intentionally small. The Streamlit app owns the layout, the
session state, and the language toggle. The agent module owns the prompt
construction and the model call. The chart generator turns a
\texttt{chart\_spec} into a Plotly figure. The voice component is a thin
wrapper around a browser API. The data loader reads the CSV outputs of
the previous sprints into a single in-memory \texttt{DataContext} that the
agent passes to the model as schema + summary text.

\section{Implementation}

\subsection{Loading the Data Context}

\texttt{chatbot/data\_loader.py} reads eleven CSV files into a shared
\texttt{DataContext} dataclass:

\begin{itemize}[noitemsep]
  \item From \texttt{output/powerbi\_star/}: \texttt{fact\_transaction},
        \texttt{fact\_session}, \texttt{dim\_member}.
  \item From \texttt{output/forecasts/}: \texttt{forecast\_revenue},
        \texttt{session\_volume}, \texttt{member\_segments},
        \texttt{member\_loyalty}, \texttt{anomalies},
        \texttt{peak\_hours\_by\_hour}, \texttt{peak\_hours\_by\_day},
        \texttt{stock\_replenishment}.
\end{itemize}

For each table the loader emits a short schema line (column names plus
dtypes) and aggregates a one-shot summary string: total revenue, member
count, loyalty-tier distribution, the next four weekly and three monthly
revenue forecasts, recent-anomaly count, average session duration, and
session-type breakdown. The two strings (\texttt{schemas} and
\texttt{summary}) are spliced into the system prompt at every turn so the
model always has both the structure of the data and the headline numbers
in front of it.

A useful side effect: the loader also exposes a
\texttt{recent\_anomaly\_count()} helper, so the Streamlit app can change
its startup greeting when something interesting has happened in the
previous seven days.

\subsection{The Conversational Agent}

\texttt{chatbot/agent.py} is the bridge between the Streamlit app and
the Gemini API. It does three things:

\begin{enumerate}
  \item \textbf{Build the system prompt.} The prompt mixes a fixed
        instruction block, the language directive, the schema string, and
        the summary string. The instructions tell the model to answer
        concisely, to include a single \texttt{chart\_spec} JSON block
        only when a chart would actually help, to suggest two or three
        follow-up questions, to refuse questions that cannot be answered
        from the available data, and to never invent numbers.
  \item \textbf{Call the model.} The full conversation history is
        re-sent on every turn (Gemini \texttt{gemini-3-flash-preview}),
        with the system prompt attached as
        \texttt{system\_instruction}. The model returns one piece of
        free text that contains the answer and, optionally, two JSON
        code blocks.
  \item \textbf{Parse the response.} A small regex pass extracts the
        \texttt{chart\_spec} block and the \texttt{suggestions} block,
        then strips them out of the visible answer so the chat shows
        clean prose only.
\end{enumerate}

The system prompt also contains a short ``source selection rules''
section that nudges the model towards the right table for each kind of
question (revenue trend $\to$ \texttt{fact\_transaction}; future
revenue $\to$ \texttt{forecast\_revenue} with a granularity filter;
hourly pattern $\to$ \texttt{peak\_hours\_by\_hour}; and so on). This
turned out to matter more than tuning the temperature. Without the rules,
the model would happily plot \texttt{fact\_session} on a non-existent
date axis or use the \texttt{anomalies} table for generic revenue
questions.

\subsection{Chart Generation}

\texttt{chatbot/chart\_generator.py} converts a \texttt{chart\_spec} dict
into a Plotly Express figure. The spec is small on purpose --- just
\texttt{type}, \texttt{x}, \texttt{y}, \texttt{source} (the table key),
\texttt{filter} (a pandas \texttt{query} expression that the agent can
fill in), and \texttt{title}. Four chart types are supported:
\textit{bar}, \textit{line}, \textit{scatter}, and \textit{pie}. If the
\texttt{y} column is non-numeric, the generator groups by \texttt{x} and
counts instead of trying to aggregate a string column. A dark theme is
applied to match the rest of the Streamlit layout.

Charts are not rendered when the model decides one is not useful (for
example, ``how many members do I have in total?'' returns a one-line
answer with no figure). This keeps the right-hand panel meaningful
instead of being filled with redundant visuals.

\subsection{Voice Input}

\texttt{chatbot/voice\_component.py} declares a custom Streamlit
component pointing at a small HTML/JS frontend in
\texttt{chatbot/voice\_html/}. The frontend uses the browser's
\texttt{SpeechRecognition} API to capture audio, transcribe it locally
in the browser, and post the transcript back to Streamlit. The selected
language toggle is forwarded to the component as a recognition locale
(\texttt{en-US}, \texttt{fr-FR}, or \texttt{ar-TN}), so French users
get French transcription and Arabic users get Tunisian-Arabic
transcription out of the box.

Keeping the recognition in the browser avoids having to wire a separate
speech-to-text service, and means voice mode works offline once the page
is loaded.

\subsection{Multilingual UI}

The \texttt{STRINGS} dictionary in \texttt{app.py} carries every visible
label in three languages. Switching language flips the dict, sends a new
language directive into the system prompt, and updates the voice
component's recognition locale --- all in one place, all without
reloading the data. The Arabic strings are stored as-is, and Streamlit
renders them right-to-left where appropriate.

The greeting is also context-aware. If the anomaly detector flagged
something in the previous seven days, the assistant opens with
``I noticed \texttt{N} anomaly(ies) in the past 7 days. Want me to show
you?'' instead of the generic ``Hi! Ask me anything.'' This tiny detail
was the change that made the assistant feel most useful in practice:
the owner does not have to think to discover that something is off.

\subsection{Quick Actions}

Four shortcut buttons sit under the chat history --- \textit{Full
summary}, \textit{Give me advice}, \textit{Any anomalies?}, and
\textit{Forecast}. Each one just enqueues a pre-written prompt into
\texttt{st.session\_state.pending\_query} and triggers a rerun, so the
model receives the question as if the user had typed it. They cover the
three or four questions the owner asks most often, and skip the typing
cost.

\section{Sprint 5 Review}

\subsection{Deliverables}

\begin{itemize}[noitemsep]
  \item \texttt{app.py} --- 252-line Streamlit application: language
        toggle, anomaly-aware greeting, chat panel with suggestion
        chips, chart panel, quick-action buttons, voice button.
  \item \texttt{chatbot/agent.py} --- prompt builder, Gemini call, and
        structured-response parser.
  \item \texttt{chatbot/chart\_generator.py} --- Plotly-Express
        renderer for the four supported chart types.
  \item \texttt{chatbot/data\_loader.py} --- shared
        \texttt{DataContext}, schema/summary builder,
        \texttt{recent\_anomaly\_count()} helper.
  \item \texttt{chatbot/voice\_component.py} +
        \texttt{chatbot/voice\_html/} --- browser-side speech
        recognition wired into Streamlit.
\end{itemize}

\subsection{Validation}

The assistant was validated against three kinds of questions:

\begin{itemize}
  \item \textbf{Descriptive} --- ``Which cashier had the best March?'',
        ``What is the average session duration?'', ``Show me revenue by
        payment method.'' Answers and charts matched the corresponding
        Power~BI page on every test case.
  \item \textbf{Forecast} --- ``How much revenue should I expect next
        month?'', ``Will members grow in May?'' Returned the values
        from \texttt{forecast\_revenue.csv} and
        \texttt{session\_volume.csv} with the right granularity filter.
  \item \textbf{Anomaly} --- ``Anything unusual last week?'' Returned
        the rows from \texttt{anomalies.csv} flagged within the
        seven-day window, matching the \texttt{recent\_anomaly\_count}
        helper.
\end{itemize}

The multilingual handling was checked end-to-end: a French question
produced a French answer with the chart title in French, and an Arabic
question produced an Arabic answer rendered right-to-left.

\section{Conclusion}

Sprint~5 turned a folder of CSV files into something the owner can
actually \emph{converse} with. There is no new data in the system ---
the agent reads exactly the same outputs the dashboards do --- but the
interaction model is fundamentally different. Instead of clicking
through four dashboard pages to answer a question, the owner asks the
question. Instead of remembering which CSV holds the anomaly list, the
greeting brings the answer to him. The chatbot is the practical face
of every previous sprint: without the cleaning, the schema, the
dashboards, and the ML layer it would have nothing to talk about, but
without the chatbot none of those would have been accessible without
opening Power~BI. The next chapter wraps up the project as a whole.
